\section{Jan 9} 

  Assume knowledge of 403 and 531 basically. Dunford \& Schwartz. Recall the definition of a \textit{normed vector space}. It induces a metric and hence a topology. A metric space is a Banach space if it is complete. It's often that if you have infinite-dim vector spaces, you have a lot of topologies to choose from, each of which have different roles. We will find out how to utilize each of them. Some examples are $\mathbb{R}^n$ with the $p$-norm, $\ell^p$ such that $\big( \sum_{j=1}^\infty |x_j|^p \big) < \infty$, or $L^p (\mathbb{R}^n)$. We will begin with the simplest topology. 

  \begin{definition}[Bounded Linear Operators]
    Let $(X, \|\cdot\|_X)$ and $(Y, \|\cdot\|_Y)$ be 2 normed spaces. Then, a linear map $A: X \to Y$ is called \textbf{bounded} if there exists a $c \geq 0$, s.t. $\forall v \in X$, 
    \begin{equation}
      \|Av \|_Y \leq c \|v\|_X
    \end{equation}
    The smallest such $c$ is denoted $\|A\| = \|A\|_{\mathrm{sup}} = \|A\|_\infty$, i.e. the spectral norm. i.e. 
    \begin{equation}
      \|A\| \coloneqq \sup_{w \in X \setminus \{0\}} \frac{\|Aw\|_Y}{\|w\|_X}
    \end{equation}
    Note that due to linearity, we can take the sup over the unit sphere as well. 
  \end{definition}

  Set 
  \begin{equation}
    L(X, Y) \coloneqq \{A : X \to Y \coloneq A \text{ linear, bounded}\}
  \end{equation}
  which is a normed vector space. The topology is called the \textit{uniform operator topology}. 

  \begin{theorem}[1.2.2.]
    Let $X, Y$ be two normed vector spaces. Let $A: X \to Y$ be linear. TFAE. 
    \begin{enumerate}
      \item $A$ is bounded. 
      \item $A$ is continuous. 
      \item $A$ is continuous at $0$ (or any point). 
    \end{enumerate}
  \end{theorem}
  \begin{proof}
    Listed. 
    \begin{enumerate}
      \item $(1 \to 2)$. The one-liner is that $A$ is Lipschitz, which is continuous, since 
        \begin{equation}
          \| Au - Av \|_Y = \| A(u - v) \|_Y \leq \|A\| \|u - v\|
        \end{equation}
      \item $(2 \to 3)$. Clear. 
      \item $(3 \to 1)$. For every $\epsilon > 0$, $\exists \delta_\epsilon > 0$ s.t. 
        \begin{equation}
          A^{-1} \big( B_Y (0, \epsilon) \big) \supset B_X (0, \delta_\epsilon)
        \end{equation}
        where $A^{-1}$ indicates the preimage. In other words, $|u|_X < \delta_\epsilon \implies |Au|_Y < \epsilon$. Let $w \in X \setminus \{0\}$. Then, 
        \begin{equation}
          w = |w| \cdot \frac{2}{\delta_\epsilon} \frac{w}{|w|} \frac{\delta_\epsilon}{2}
        \end{equation}
        So by linearity, we have 
        \begin{equation}
          Aw = \frac{|w| 2}{\delta_\epsilon} A \bigg( \frac{w}{|w|} \frac{\delta_\epsilon}{2} \bigg) \implies |A w| \leq \frac{2 |w|}{\delta_\epsilon} \epsilon
        \end{equation}
        So 
        \begin{equation}
          \|A\| \leq \frac{2 \epsilon}{\delta_\epsilon}
        \end{equation}
    \end{enumerate}
  \end{proof}

  So basically, boundedness is equivalent to continuity, which is why most analysts look at bounded operators. Let's review some notation from linear algebra, given $A \in L(X, Y)$. Then, 
  \begin{enumerate}
    \item $\mathrm{ker}(A) \subset X$ is a closed subspace. This was obvious in finite dimensions, but in infinite dimensions, you need the boundedness. 
    \item $\mathrm{Im}(A) \subset Y$ is a subspace, but it is not closed, which is not as good. 
    \item $\mathrm{koker}(A) \coloneqq Y / \mathrm{Im}(A)$ is also a vector space, but it can be nasty, since if $\mathrm{Im}(A)$ is not closed, then the quotient space is not Hausdorff. 
  \end{enumerate} 

  \begin{definition}
    2 norms $|\cdot|_1, |\cdot|_2$ on $X$ are equivalent if $\exists c_1 \leq c_2$ s.t. $\forall v \in X$, 
    \begin{equation}
      c_1 |v_1| \leq |v|_2 \leq c_2 |v_1|
    \end{equation}
  \end{definition}

  It is easy to see that equivalent norms induce the same topology, and equivalence of norms is an equivalence relation. 

  \begin{theorem}
    Let $X$ be any finite-dimensional vector space. Then, any two norms on $X$ are equivalent. 
  \end{theorem}
  \begin{proof}
    Choose a basis $\{e_1, \ldots, e_n\}$ for $X$. Let's define the Euclidean norm $\| \sum_{j=1}^n x_j e_j\|_2 \coloneqq \sqrt{\sum_j x_j^2}$. Let $|\cdot|$ be any other norm on $X$. Then, by the triangle inequality, followed by Cauchy-Schwartz, 
    \begin{equation}
      \bigg| \sum_{j=1}^n x_j e_j \bigg| \sum_j |x_j| |e_j|  \leq \sqrt{\sum_j x_j^2} \underbrace{\sqrt{\sum_j |e_j|^2}}_{= c_1} \implies |v| \leq c_1 \|v\|_2
    \end{equation}

    Now let $S^{n-1} = \|\cdot\|^{-1} (1)$, i.e. the preimage of the Euclidean norm of $1$. So $| \cdot |$ is continuous w.r.t. the $\|\cdot\|_2$-topology,\footnote{it is trivially continuous w.r.t. its own topology that it induces.} Since it is well known that $S^{n-1}$ is compact, $|\cdot|$ attains its minimum on $S^{n-1}$, call it $c_2$. For $v \neq 0$, 
    \begin{equation}
      |v| = \|v\|_2 \bigg| \frac{v}{\|v\|_2} \bigg| \geq c_2 \|v_2\|
    \end{equation}

    We mixed up the $c_1, c_2$ labels, but the idea is the same. 
    \begin{equation}
      c_2 \|v\|+2 \leq |v| \leq c_1 \|v\|_2 
    \end{equation}
  \end{proof}

  This establishes that we know the metric structure of $\mathbb{R}^n$ very well, and so $\mathbb{R}^n$ is pretty boring. Unless you're in stats and you study very high-dimensional analysis. 

  \begin{corollary}[1.2.6]
    Every finite dimensional normed vector space is Banach. 
  \end{corollary}
  \begin{proof}
    Every finite-dimensional vector space is isomorphic to $\mathbb{R}^n$, which is complete under all norms. 
  \end{proof}

  \begin{corollary}[1.2.7]
    Let $X$ be a normed vector space. Then every finite-dimensional subspace is closed. 
  \end{corollary}
  \begin{proof}
    Let $Y \subset X$ be finite dimensional. Then $Y$ is complete, hence is closed. 
  \end{proof}

  \begin{corollary}
    Let $X$ be a finite-dimensional normed vector space. Then $K \subset X$ is compact iff it is closed and bounded. 
  \end{corollary}

  \begin{corollary}[1.2.9]
    Let $X, Y$ be two normed vector spaces, with $\dim(X) < +\infty$. Then every linear function $f: X \to Y$ is bounded. 
  \end{corollary}
  \begin{proof}
    Let $A: X \to Y$ be linear. Define the graph norm on $X$ to be 
    \begin{equation}
      |w|_A \coloneqq |w|_X + |Aw|_Y
    \end{equation}
    But this is equivalent to Euclidean norm since all norms on finite-dimensional vector spaces are equivalent. So $\exists c_1$ s.t. 
    \begin{equation}
      |w|_X + |Aw|_Y \leq c_1 |w|_X \implies |Aw|_Y \leq c_1 |w|_X
    \end{equation}
  \end{proof}

  So it tells us that any linear map from a finite-dimensional vector space to any vector space is continuous. 

  \begin{theorem}[1.2.11]
    Let $X$ be a normed vector space, $B \subset X$ denote the unit ball, $S \subset X$ denote the unit sphere, then TFAE. 
    \begin{enumerate}
      \item $\dim(X) < +\infty$. 
      \item $B$ is compact. 
      \item $S$ is compact. 
    \end{enumerate}
  \end{theorem}
	\begin{proof}
		If this is Hilbert space, then the proof is trivial. But this is hard in Banach spaces. It's analogous to doing linear vs nonlinear PDEs. 
	\end{proof}

\section{Jan 16} 

	You should be very careful with your intuition in an infinite-dimensional vector space. You should develop this intuition on what is or isn't obvious in an infinite-dimensional vector space, so we're going over some seemingly ``obvious'' results. 

	\begin{definition}[Hilbert Space]
		A \textbf{Hilbert space} is a complete inner product space. 
	\end{definition}

	In analysis, you want to solve problems by construting sequences. It's bad if your vector space is not complete, i.e. these sequences don't converge. 

	\begin{lemma} 
	  Let $K$ be a closed, nonempty convex subset of Hilbert space $H$. Then, there exists a unique element of minimal norm in $K$. 
	\end{lemma}
	\begin{proof}
		You should remember this proof since it's used in the Riesz Representation theorem's proof. Let $\delta = \inf_{x \in K} |x|$. Pick a sequence $(x_j) \subset K$ s.t. $|x_j| \to \delta$ as $j \to +\infty$. We claim that $(x_j)$ is Cauchy. This is somewhat nonobvious, so as a proof, let $\epsilon > 0$. Then, there exists a $N \in \mathbb{N}$ s.t. $j > N \implies |x_j|^2 \leq \delta^2 + \frac{\epsilon^2}{4}$. So by convexity, 
		\begin{equation}
			m, n > N \implies \frac{x_m + x_n}{2} \in K
		\end{equation}
		Since $\delta$ is an infimum, 
		\begin{equation}
			\bigg| \frac{x_m + x_n}{2} \bigg| \geq \delta \implies |x_m + x_n|^2 \geq 4 \delta^2
		\end{equation}
		By expanding the quadratic 
		\begin{align}
			|x_m - x_n|^2 & = |x_m|^2 + |x_n|^2 - 2 \langle x_m, x_n \rangle \\ 
										& = - |x_m + x_n|^2 + 2|x_m|^2 + 2|x_n|^2 \\ 
										& < - 4 \delta^2 + 4 \delta^2 + 4 \frac{\epsilon^2}{4} = \epsilon^2
		\end{align} 
		which implies that $|x_m - x_n| < \epsilon$ if $m, n > N$. So, it is Cauchy\footnote{Check if this is related to strongly Cauchy?} and since $H$ is complete, $(x_j)$ has a limit---call it $x_\infty \in K$ (since $K$ is closed) with $|x_\infty| = \delta$. Now we prove uniqueness. Suppose $y \in K$ satisfies $|y| = \delta \implies \frac{x_\infty + y}{2} \in K \implies |x_\infty + y|^2 > 4 \delta^2$. So, $|x_\infty - y|^2 = 0 \implies x_\infty = y$. 
	\end{proof}
    

	\begin{theorem}[Riesz Representation Theorem]
	  Let $H$ be a Hilbert space. Define $\Lambda: H \to H^\ast$ by $\Lambda(x)(y) \coloneqq \langle y, x \rangle$. Then, $\Lambda$ is an isometeric isomorphism. 
	\end{theorem}
	\begin{proof}
	  For $x \neq 0$, we have 
		\begin{equation}
			\|\Lambda (x)\| = \sup_{y \neq 0} \frac{|\Lambda(x)(y)|}{|y|} = \sup_{y \neq 0} \frac{|\langle y, x \rangle|}{|y|} \leq \frac{|x||y|} {|y|} = |x|
		\end{equation}
		For the other direction, 
		\begin{equation}
			\frac{|\Lambda(x)(x)|}{|x|} = \frac{|x|^2}{|x|} = |x|
		\end{equation}
		Therefore, $\|\Lambda(x)(x)\| = |x|$, and $\Lambda$ is an isometry. 
		\begin{enumerate}
		  \item \textit{Injective}. It is easy to show that $\Lambda: H \to H^\ast$ is injective. 
			\item \textit{Surjective}. This is more involved, and we need convexity. For simplicity, let our field be $\mathbb{R}$. We've shown that $\Lambda: H \to H^\ast$ is an isometric injection. Let $\phi \in H^\ast \setminus \{0\}$. There are 2 families of natural convex sets: affine spaces and balls. It is easy to see that $K = \phi^{-1} (1)$ is an affine space and hence convex. Let $x_0$ be an element of minimal norm of $K$. We claim that $x_0 \perp \mathrm{ker}(\phi)$. Indeed, we have 
				\begin{align}
					y \in \mathrm{ker}(\phi) & \implies x_0 + ty \in K \text{ for all } t \in \mathbb{R} \\ 
																	 & \implies |x_0|^2 \leq |x_0 + ty|^2 = |x_0|^2 + 2t \langle x_0, y \rangle + t^2 |y|^2 \\ 
																	 & \implies 0 \leq 2t \langle x_0, y \rangle + t^2 |y|^2
				\end{align}
				Take $t = -\frac{\langle x_0, y \rangle}{|y|^2}$, then we have 
				\begin{equation}
					0 \leq \frac{-2 \langle x_0, y \rangle^2}{|y|^2} + \frac{\langle x_0, y \rangle^2}{|y|^2} \implies 0 \leq - \frac{\langle x_0, y \rangle^2}{|y|^2} 
				\end{equation}
				and so $\langle x_0, y \rangle = 0$. So to summarize, $\phi(x_0) = 1$ with $x_0 \perp \mathrm{ker}(\phi)$. Let $x \in H$. Then, 
				\begin{equation}
					\phi\big( x - \phi(x) x_0 \big) = \phi(x) - \phi(x) \underbrace{\phi(x_0)}_{=1} = 0
				\end{equation}
				But this means 
				\begin{align}
					x - \phi (x) x_0 \in \mathrm{ker}(\phi) 
						& \implies 0 = \langle x - \phi(x) x_0, x_0 \rangle \\ 
						& \implies \langle x, x_0 \rangle = \phi(x) \langle x_0, x_0 \rangle = \phi(x) |x_0|^2 \\ 
						& \implies \phi(x) = \bigg\langle x, \frac{x_0}{|x_0|^2} \bigg\rangle \\ 
						& \implies \phi = \Lambda \bigg( \frac{x_0}{|x_0|^2} \bigg)
				\end{align}
				and therefore, $\Lambda$ is surjective. 
		\end{enumerate} 
	\end{proof}

	\begin{definition}[Orthogonal Subspace]
	  Let $S \subset H$ be a subspace of a Hilbert space, and let 
		\begin{equation}
			S^\perp \coloneqq \{x \in H \colon \langle s, x \rangle = 0 \forall s \in S\}
		\end{equation} 
	\end{definition}


	\begin{theorem}[1.4.6]
	  Let $E \subset H$ be a closed subspace of Hilbert space $H$. Then, $H = E \oplus E^\perp$. 
	\end{theorem}
	\begin{proof}
	  If $x \in E \cap E^\perp$, then $x \perp x \iff \langle x, x = 0 \iff x = 0$. If $x \in H$, then $x + E$ is a closed convex set, which implies that there exists an element $x_0$ of least norm in $x + E$. Therefore, $x_0 \perp E$ as before (previous argument: $|x_0 + ty|^2 \geq |x_0|^2$...), which implies that 
		\begin{equation}
			x = \underbrace{x - x_0}_{\in E} + \underbrace{x_0}_{\in E^\perp}
		\end{equation}
	\end{proof}

	More elementary properties are in the book, but now about Banach algebras, which are infinite-dimensional analogue of matrices. Remember $\mathcal{L}(X) \coloneqq \mathcal{L}(X, X)$ as the set of bounded endomorphisms on $X$. Remember if $Y$ is Banach, then $L(X, Y)$ is Banach w.r.t. the operator norm. 

	\begin{definition}[Banach Algebra]
	  A \textbf{Banach algebra} is a Banach space $A$ equipped with a bilinear map 
		\begin{equation}
		  A \times A \to A, \qquad (a, b) \to ab
		\end{equation}
		called the \textit{product}, satisfying for all $a, b, c \in A$, 
		\begin{enumerate}
		  \item \textit{Associativity}. $(ab)c = a(bc)$. 
			\item \textit{Respects norm}. $|ab| \leq |a| |b|$. 
		\end{enumerate}
		It is called 
		\begin{enumerate}
		  \item \textbf{commutative} if $ab = ba$ for all $a, b \in A$. 
			\item \textbf{unital} if there exists a $1_A \in A$ s.t. $1_A a = a 1_A = a$ for all $a \in A$. 
		\end{enumerate}
	\end{definition}

	\begin{definition}
		An element $a \in A$ of a unital Banach algebra is called \textbf{invertible} if there exists a $b \in A$ s.t. $ab = ba = 1_A$. It is unique if it exists and is denoted $a^{-1}$. 
	\end{definition}

	Here is the basic example to hold in mind. 

	\begin{theorem}[Banach Algebra of Bounded Linear Endomorphisms]
		$\mathcal{L}(X)$ is a unital Banach algebra w.r.t composition as multiplication. 
	\end{theorem}
	\begin{proof}
	  The identity map is the unity element. It is associate since composition is associative. So, we just need to prove the norm inequality. Define 
		\begin{equation}
			|ab| \coloneqq \sup_{x \neq 0} \frac{|abx|}{|x|} = \sup_{x \neq 0} \frac{|a(bx)|}{|x|} \leq \sup_{x \neq 0} \frac{|a| |bx|}{|x|} = |a| \cdot \sup_{x \neq 0} \frac{|bx|}{|x|} = |a| |b|
		\end{equation}
	\end{proof}

	Note that this inequality can be strict since 
	\begin{equation}
		\underbrace{\begin{pmatrix} 1 & 0 \\ 0 & 0 \end{pmatrix}}_{1} \underbrace{\begin{pmatrix} 1 & 0 \\ 0 & 0 \end{pmatrix}}_{1} = \underbrace{\begin{pmatrix} 1 & 0 \\ 0 & 0 \end{pmatrix}}_{0}
	\end{equation}
	or just take any invertible map $A, A^{-1}$ where $A A^{-1} = I$. 

	\begin{definition}[Compact Linear Operator]
		A $f \in \mathcal{L}(X)$ (or $\mathcal{L}(X, y)$) is called \textbf{compact} if for all bounded sets $B \subset X$, $f(B)$ is precompact, i.e. $\overline{f(B)}$ is compact. So, define 
		\begin{equation}
			K(X) \coloneqq \{ f \in \mathcal{L}(X) \text{ s.t. } f \text{ is compact}\}
		\end{equation}
	\end{definition}

	$K(X)$ is very analogous to finite-rank matrices and are the first thing to study. 

	\begin{theorem}[Compact Linear Endomorphisms are Banach Algebra]
	  $K(X)$ is a Banach algebra, but it is nonunital since the identity is not a compact map. 
	\end{theorem}

	The following example comes up a lot in K-theory, analysis, and most of algebraic geometry is based on this. 

	\begin{example}
	  Let $\Omega$ be a metric space. Then, $C(\Omega)$ is a commutative unital Banach algebra w.r.t. the supremum norm. 
	\end{example}
